\documentclass[11pt]{article}
\usepackage[top=1.00in, bottom=1.0in, left=1in, right=1in]{geometry}
\renewcommand{\baselinestretch}{1.1}
\usepackage{graphicx}
\usepackage{natbib}
\usepackage{amsmath}
\usepackage{gensymb}
\usepackage{parskip}
\usepackage{xcolor}
\usepackage[normalem]{ulem}

\def\labelitemi{--}
\parindent=0pt

\begin{document}
\renewcommand{\refname}{\CHead{}}


\title{Supplement: The problem and promise of \\ modeling plant leafout under warmer winters} %  How overconfident models of tree leafout have stalled progress... Infinite models of leafout given warming winters...? % I don't love this title.
\author{E.M. Wolkovich$^1$, Justin Ngo$^1$, Victor van der Meersch$^1$, Jonathan Auerbach$^2$} % also possibly: Rob Guy -- I emailed him in mid-December 2025, have not heard back
\date{\today}
\maketitle

$^1$ Forest \& Conservation Sciences, Faculty of Forestry, University of British Columbia, 2424 Main Mall, Vancouver, BC V6T 1Z4, Canada\\
$^2$ Department of Statistics, George Mason University, 4511 Patriot Circle, Fairfax, VA 22030, United States\\

\emph{The search for better indicators}

Increasing evidence suggests that the polysaccharide callose and its degradation through 1,3-$\beta$-glucanases plays a pivotal role in bud endodormancy \citep{rinne2011,vanderschoot2014,pan2023epigenetic}, potentially alongside other factors, such as phytohormones \citep[][]{tylewicz2018photoperiodic,pan2021aba}, though whether any of these will prove as a useful measurable indicator is still unclear. Currently, measuring callose deposition and degradation may be possible through upregulated expression of 1,3-$\beta$-glucanases  localized to the plasmodesmata, though this is far from simple. Similarly, although ABA and gibberellin are implicated, they are not easy to measure. Finding a particular transcript or ratio of transcripts would be likely far easier to measure that callose-related compound, ABA or gibberellin. 

% START HERE ... WORK on adding in what Rob sent ...

% From Justin: 
% Researcher could measure upregulation expression of 1,3-$\beta$-glucanases localized to the plasmodesmata during the overwintering/endodormancy period, providing a method of quantifying and visualizing callose deposition and degradation in the bud during chilling. Likewise, phytohormones heavily involved in dormancy induction and release like ABA and GA have similarly measurable counterparts, such as the ABA-mediated expressions of [MADS-box associated proteins/morphogenesis and development associated proteins/ PpDAM6] and [MAP kinases/protein activation and deactivation cascade-associated enzymes/PpMAPK6] in peaches.

% From Rob Guy (Jan 2026): 
% Yes, callose or the glucanases may be appropriate indicators. Ideally, though, it should be something easily measureable, like a particular transcript or ratio of transcripts (perhaps for the glucanase). Although ABA and gibberellin are implicated, they are not easy to measure. I'll attach some text I wrote for another purpose (but never used) if you want to flesh this out a bit.


\bibliographystyle{/Users/Lizzie/Documents/EndnoteRelated/Bibtex/styles/besjournals}
\section{References}
\bibliography{..//refs/bibforchillingconceptms}

\end{document}
